%% Appendix: supplementary results (MedIA revision, v2 pipeline).
%% Figures: paper/media/figs (rendered from snapshot 1381bdd759f9 by src/v2/make_figures_v2.py --freeze).
%% Tables: paper/media/tables (src/v2/paper_tables.py). Ladder-only rank agreement:
%% results/v2/analysis/S_ladder_only/cross/kendall_id.csv.
\section{Compact-tier results}
\label{app:compact}

This section gives the compact-tier counterparts of the main-text figures
and tables. Table~\ref{tab:main-C} lists worst-domain and mean AUROC,
Fig.~\ref{fig:perdomain-C} the per-domain results, Fig.~\ref{fig:variance-C}
the variance decomposition, Fig.~\ref{fig:power-C} the power curves and
Fig.~\ref{fig:ranks-C} the rankings. Figure~\ref{fig:seedspread} shows the
individual seeds of three compact-tier arms next to the best standard-tier arm.

\begin{table}[width=\linewidth,pos=p]
\caption{Compact tier: worst-domain AUROC (rank among the 20 arms) and mean
AUROC over held-out domains, ID-validation selection, five seeds. Bold: best
worst-domain AUROC in the cohort. $^{\dagger}$Transductive.}
\label{tab:main-C}
\footnotesize
\setlength{\tabcolsep}{3pt}
% generated by src/v2/paper_tables.py from results/v2/analysis/C/*/summary_id.csv, ranks_id.csv
\begin{tabular*}{\tblwidth}{@{\extracolsep{\fill}}l*{4}{rr}@{}}
\toprule
Arm & \multicolumn{2}{c}{Camelyon17 ($K=5$)} & \multicolumn{2}{c}{Canine cSCC ($K=5$)} & \multicolumn{2}{c}{MIDOG 2021 ($K=3$)} & \multicolumn{2}{c}{MIDOG++ ($K=7$)} \\
\cmidrule(lr){2-3}\cmidrule(lr){4-5}\cmidrule(lr){6-7}\cmidrule(lr){8-9}
 & worst (rank) & mean & worst (rank) & mean & worst (rank) & mean & worst (rank) & mean \\
\midrule
\multicolumn{9}{@{}l}{\emph{Reference and controls}}\\
ERM & 0.698 (13) & 0.844 & 0.551 (18) & 0.714 & 0.765 (17) & 0.801 & 0.766 (13) & 0.827 \\
Compute-matched ERM & 0.727 (12) & 0.848 & 0.466 (20) & 0.693 & 0.789 (13) & 0.818 & \textbf{0.862} (1) & 0.887 \\
\addlinespace[2pt]
\multicolumn{9}{@{}l}{\emph{Domain-aware training}}\\
GroupDRO & 0.670 (14) & 0.830 & 0.579 (17) & 0.715 & 0.764 (18) & 0.802 & 0.731 (16) & 0.804 \\
IRM & 0.607 (19) & 0.828 & 0.589 (16) & 0.718 & 0.785 (14) & 0.809 & 0.741 (15) & 0.805 \\
CORAL & 0.637 (17) & 0.816 & 0.596 (15) & 0.723 & 0.803 (10) & 0.808 & 0.802 (9) & 0.845 \\
MixStyle & 0.665 (15) & 0.849 & 0.605 (13) & 0.735 & 0.777 (15) & 0.794 & 0.709 (17) & 0.794 \\
Fish & 0.615 (18) & 0.842 & 0.502 (19) & 0.676 & 0.676 (20) & 0.759 & 0.667 (20) & 0.765 \\
LISA & 0.735 (11) & 0.846 & 0.669 (10) & 0.718 & 0.714 (19) & 0.762 & 0.691 (18) & 0.778 \\
\addlinespace[2pt]
\multicolumn{9}{@{}l}{\emph{Stain normalization and color}}\\
Macenko$^{\dagger}$ & 0.904 (4) & 0.943 & 0.698 (7) & 0.813 & 0.870 (2) & 0.878 & 0.828 (2) & 0.868 \\
Macenko (per tile) & 0.524 (20) & 0.861 & 0.671 (9) & 0.763 & 0.844 (7) & 0.861 & 0.816 (3) & 0.849 \\
H\&E jitter & 0.904 (3) & 0.933 & 0.674 (8) & 0.795 & 0.864 (5) & 0.876 & 0.809 (8) & 0.843 \\
TIA-style$^{\dagger}$ & 0.889 (6) & 0.951 & 0.730 (6) & 0.777 & 0.864 (6) & 0.879 & 0.810 (7) & 0.853 \\
Aug-only control & 0.815 (10) & 0.921 & 0.745 (5) & 0.766 & 0.805 (9) & 0.819 & 0.670 (19) & 0.745 \\
\addlinespace[2pt]
\multicolumn{9}{@{}l}{\emph{Self-supervision}}\\
SimCLR & 0.897 (5) & 0.939 & 0.613 (12) & 0.734 & 0.795 (12) & 0.806 & 0.789 (10) & 0.838 \\
SimCLR+SAM & \textbf{0.909} (1) & 0.944 & 0.603 (14) & 0.738 & 0.797 (11) & 0.801 & 0.781 (12) & 0.830 \\
\addlinespace[2pt]
\multicolumn{9}{@{}l}{\emph{Sharpness-aware}}\\
SAM & 0.639 (16) & 0.825 & 0.614 (11) & 0.725 & 0.776 (16) & 0.809 & 0.757 (14) & 0.815 \\
\addlinespace[2pt]
\multicolumn{9}{@{}l}{\emph{Test-time adaptation}}\\
BN-adapt (ERM)$^{\dagger}$ & 0.854 (9) & 0.920 & 0.801 (2) & 0.836 & 0.865 (3) & 0.880 & 0.814 (4) & 0.834 \\
TENT (ERM)$^{\dagger}$ & 0.874 (7) & 0.933 & 0.799 (3) & 0.835 & 0.865 (4) & 0.878 & 0.814 (5) & 0.830 \\
BN-adapt (SimCLR)$^{\dagger}$ & 0.858 (8) & 0.938 & \textbf{0.822} (1) & 0.836 & 0.842 (8) & 0.852 & 0.786 (11) & 0.830 \\
BN-adapt (H\&E jitter)$^{\dagger}$ & 0.908 (2) & 0.950 & 0.798 (4) & 0.831 & \textbf{0.877} (1) & 0.891 & 0.812 (6) & 0.845 \\
\bottomrule
\end{tabular*}

\end{table}

\begin{figure}[pos=p]
  \centering
  \includegraphics[width=\linewidth]{fig_perdomain_C}
  \caption{Test AUROC of every compact-tier arm at every held-out domain (mean
  over five seeds; ID-validation selection). Outlined: each arm's worst domain.
  $^{\dagger}$Transductive.}
  \label{fig:perdomain-C}
\end{figure}

\begin{figure}[pos=p]
  \centering
  \includegraphics[width=0.92\linewidth]{fig_variance_C}
  \caption{Variance decomposition at the compact tier, as in
  Fig.~\ref{M-fig:variance}. $^{\dagger}$Transductive.}
  \label{fig:variance-C}
\end{figure}

\begin{figure}[pos=tbp]
  \centering
  \includegraphics[width=\linewidth]{fig_power_C}
  \caption{Domains needed for 80\% power at the compact tier, as in
  Fig.~\ref{M-fig:power}.}
  \label{fig:power-C}
\end{figure}

\begin{figure}[pos=p]
  \centering
  \includegraphics[width=0.85\linewidth]{fig_ranks_C}
  \caption{Rankings at the compact tier, as in Fig.~\ref{M-fig:ranks}.
  $^{\dagger}$Transductive.}
  \label{fig:ranks-C}
\end{figure}

\begin{figure}[pos=tbp]
  \centering
  \includegraphics[width=\linewidth]{fig_seedspread}
  \caption{Individual seeds at every held-out domain for compact-tier ERM,
  SimCLR and H\&E jitter, and for the best standard-tier arm (fine-tuned Lunit
  DINO with H\&E jitter, the arm with the highest mean worst-domain AUROC over
  the four cohorts, 0.895). Bars mark seed means. The $y$-axes are scaled per
  cohort.}
  \label{fig:seedspread}
\end{figure}

\clearpage
\section{Model-selection rules}
\label{app:selection}

Figures~\ref{fig:selection-S} and \ref{fig:selection-C} show each arm's mean
difference from ERM under the four selection rules. All rules use the same
ten checkpoints of the same runs (Section~\ref{M-sec:protocol-lodo}), so the
differences come from selection alone. The $\tau_b$ values printed in the
figures compare the rankings of these mean differences; those of
Table~\ref{M-tab:selection} compare worst-domain rankings. At the standard tier
the rules agree closely. At the compact tier the OOD-validation and oracle
rules reorder the canine and Camelyon17 arms considerably.

\begin{figure}[pos=p]
  \centering
  \includegraphics[width=\linewidth]{fig_selection_S}
  \caption{Standard tier: mean AUROC difference from ERM under ID-validation
  (primary, with 95\% interval), OOD-validation, last-checkpoint and oracle
  selection. Each panel shows the arms analyzed under every rule available in
  the cohort. $^{\dagger}$Transductive.}
  \label{fig:selection-S}
\end{figure}

\begin{figure}[pos=p]
  \centering
  \includegraphics[width=\linewidth]{fig_selection_C}
  \caption{Compact tier: mean AUROC difference from ERM under the four
  selection rules, as in Fig.~\ref{fig:selection-S}.
  $^{\dagger}$Transductive.}
  \label{fig:selection-C}
\end{figure}

\clearpage
\section{Calibration}
\label{app:calibration}

Figures~\ref{fig:calibration-S} and \ref{fig:calibration-C} give the expected
calibration error (Eq.~\ref{M-eq:ece}) of every arm, averaged over domains and at
its worst domain, and reliability diagrams for the fold in which ERM performs
worst. At the standard tier, one of the two fine-tuned foundation models with
H\&E jitter is the best calibrated arm on every cohort. ERM is poorly calibrated on the
canine cohort and MIDOG~2021 (mean ECE 0.180 on both). Calibration varies
between domains as much as discrimination does: for most arms the worst-domain
ECE is well above the mean.

\begin{figure}[pos=p]
  \centering
  \includegraphics[width=0.92\linewidth]{fig_calibration_S}
  \caption{Calibration at the standard tier. (a) Test ECE (15 equal-width bins)
  at the worst domain (filled) and averaged over domains (open). (b) Reliability
  diagrams for the fold whose test domain is ERM's worst, pooled over seeds.
  $^{\dagger}$Transductive.}
  \label{fig:calibration-S}
\end{figure}

\begin{figure}[pos=p]
  \centering
  \includegraphics[width=0.92\linewidth]{fig_calibration_C}
  \caption{Calibration at the compact tier, as in Fig.~\ref{fig:calibration-S}.
  $^{\dagger}$Transductive.}
  \label{fig:calibration-C}
\end{figure}

\clearpage
\section{Representation ladder}
\label{app:ladder}

Figure~\ref{fig:ladder} compares nine frozen backbones under one linear-probe
protocol: four ImageNet-1k models (ResNet-18, ResNet-50, ConvNeXt-T
\citep{liu2022convnext}, ViT-B/16), a ViT-B/16 pre-trained with weak supervision on
3.6 billion images and fine-tuned on ImageNet-1k \citep{singh2022swag}, and four Lunit pathology models
\citep{kang2023benchmarking}. The Lunit Barlow Twins ResNet-50 and DINO ViT-S/16
are the best frozen encoders on the canine cohort and, by mean AUROC, on
Camelyon17; by worst-domain AUROC on Camelyon17 the DINO ViT is fourth, behind
Barlow Twins and two ImageNet models. On the mitosis cohorts every
frozen probe is well below fine-tuned ERM, and the Lunit SwAV ResNet-50 is
below the ImageNet ResNet-50 averaged over the four cohorts. Within this
probe-only ladder, the two mitosis cohorts rank the backbones almost
identically ($\tau_b=0.83$), while the two tumor cohorts barely agree
($\tau_b=0.17$).

\begin{figure}[pos=tbp]
  \centering
  \includegraphics[width=\linewidth]{fig_ladder}
  \caption{Frozen backbones with a linear probe (standard tier, 224~px input,
  three seeds, ID-validation selection): worst-domain (filled, with
  seed-bootstrap 95\% interval) and mean AUROC (open). Vertical line: fine-tuned
  ImageNet ResNet-50 ERM, worst domain.}
  \label{fig:ladder}
\end{figure}

\clearpage
\section{Sensitivity analyses}
\label{app:sensitivity}

\paragraph{Scale normalization of MIDOG~2021} Figure~\ref{fig:sensitivity}
compares the compact-tier results on the scale-normalized (primary) and the
un-normalized MIDOG~2021 cohort.

\begin{figure}[pos=tbp]
  \centering
  \includegraphics[width=0.9\linewidth]{fig_sensitivity_C}
  \caption{MIDOG~2021 with and without scale normalization (compact tier,
  20 arms, same seeds and folds): (a) worst-domain and (b) mean AUROC.
  $^{\dagger}$Transductive.}
  \label{fig:sensitivity}
\end{figure}

\paragraph{PatchCamelyon} Table~\ref{tab:pcam} gives the PatchCamelyon test
AUROC of the Camelyon17 models at the ID-validation-selected checkpoint,
averaged over seeds, with its mean and minimum over the five Camelyon17 folds.
The score is never used for selection. Test-time adaptation is not evaluated
on PatchCamelyon.

\begin{table}[width=\linewidth,pos=t]
\caption{PatchCamelyon test AUROC of the Camelyon17 models: mean and minimum
over the five folds of the seed-mean AUROC. $^{\dagger}$Transductive.}
\label{tab:pcam}
\footnotesize
% generated by src/v2/paper_tables.py from the pcam field of the frozen run records (ID-val checkpoint)
\begin{tabular*}{\tblwidth}{@{\extracolsep{\fill}}lrrrr@{}}
\toprule
 & \multicolumn{2}{c}{Compact tier} & \multicolumn{2}{c}{Standard tier} \\
\cmidrule(lr){2-3}\cmidrule(lr){4-5}
Arm & mean & min over folds & mean & min over folds \\
\midrule
ERM & 0.698 & 0.641 & 0.772 & 0.665 \\
Compute-matched ERM & 0.652 & 0.542 & -- & -- \\
GroupDRO & 0.697 & 0.634 & 0.791 & 0.738 \\
IRM & 0.716 & 0.675 & 0.753 & 0.669 \\
CORAL & 0.691 & 0.621 & 0.772 & 0.707 \\
MixStyle & 0.739 & 0.703 & 0.780 & 0.702 \\
Fish & 0.647 & 0.581 & 0.729 & 0.618 \\
LISA & 0.672 & 0.576 & 0.756 & 0.631 \\
Macenko$^{\dagger}$ & 0.765 & 0.677 & 0.752 & 0.634 \\
Macenko (per tile) & 0.756 & 0.686 & 0.837 & 0.784 \\
H\&E jitter & 0.768 & 0.664 & 0.843 & 0.801 \\
TIA-style$^{\dagger}$ & 0.802 & 0.730 & 0.858 & 0.838 \\
Aug-only control & 0.777 & 0.725 & -- & -- \\
SimCLR & 0.743 & 0.701 & 0.753 & 0.656 \\
SimCLR+SAM & 0.761 & 0.724 & -- & -- \\
SAM & 0.690 & 0.544 & 0.780 & 0.684 \\
Lunit BT FT & -- & -- & 0.838 & 0.791 \\
Lunit BT FT + H\&E jitter & -- & -- & 0.850 & 0.781 \\
Lunit DINO FT & -- & -- & 0.905 & 0.874 \\
Lunit DINO FT + H\&E jitter & -- & -- & 0.908 & 0.895 \\
RN50-IN probe & -- & -- & 0.820 & 0.806 \\
Lunit BT probe & -- & -- & 0.889 & 0.866 \\
Lunit DINO probe & -- & -- & 0.886 & 0.867 \\
\bottomrule
\end{tabular*}

\end{table}

\paragraph{Hyperparameters} Figures~\ref{fig:tuning-S} and \ref{fig:tuning-C}
show how test AUROC of the seed-0 tuning runs varies across each grid. Test
AUROC is shown only to display this sensitivity; selection used ID-validation
AUROC (Section~\ref{app:tuning}).

\begin{figure}[pos=p]
  \centering
  \includegraphics[width=0.78\linewidth]{fig_tuning_S}
  \caption{Standard-tier tuning runs (seed 0): test AUROC at the
  ID-validation-selected checkpoint for each grid value, worst fold (filled) and
  mean over folds (open). Brackets: number of folds whose ID-validation AUROC
  selected the value. `$k$ div.': $k$ folds diverged. Gray line: worst domain of
  final ERM, seed~0.}
  \label{fig:tuning-S}
\end{figure}

\begin{figure}[pos=p]
  \centering
  \includegraphics[width=0.78\linewidth]{fig_tuning_C}
  \caption{Compact-tier tuning runs (seed 0), as in Fig.~\ref{fig:tuning-S}.}
  \label{fig:tuning-C}
\end{figure}

\clearpage
\section{Response shape}
\label{app:response}

Table~\ref{tab:response} summarizes Oldham's correlation $r$ between each
arm's per-domain gain over ERM and the per-domain mean of the two
(Section~\ref{sec:protocol-response}). A negative $r$ means that the arm gains
most where both arms do worst, so it narrows the spread between domains. On the
tumor cohorts at both tiers and on MIDOG~2021 at the compact tier, most arms
have a negative $r$; on the canine cohort at the standard tier, all 20 do. The
Pitman--Morgan test is significant for 14 of them. The exact sign-flip test
cannot reach $p<0.05$ with five or fewer domains, and on MIDOG++ its smallest
$p$ is 0.078 at the standard tier. We therefore read the pattern as consistent but not established.

\begin{table}[width=\linewidth,pos=t]
\caption{Response shape: number of arms with negative Oldham's correlation $r$,
median $r$, number of arms with Pitman--Morgan $p<0.05$, and the smallest
exact sign-flip $p$ over arms.}
\label{tab:response}
\footnotesize
% generated by src/v2/paper_tables.py from response_id.csv
\begin{tabular*}{\tblwidth}{@{\extracolsep{\fill}}llrrrrr@{}}
\toprule
Tier & Cohort & $K$ & $r<0$ & median $r$ & Pitman--Morgan $p<0.05$ & smallest exact $p$ \\
\midrule
S & Camelyon17 & 5 & 15/20 & -0.36 & 7 & 0.062 \\
 & Canine cSCC & 5 & 20/20 & -0.91 & 14 & 0.062 \\
 & MIDOG 2021 & 3 & 10/20 & -0.01 & 0 & 0.250 \\
 & MIDOG++ & 7 & 4/20 & +0.13 & 5 & 0.078 \\
\addlinespace[3pt]
C & Camelyon17 & 5 & 12/19 & -0.52 & 7 & 0.125 \\
 & Canine cSCC & 5 & 17/19 & -0.73 & 6 & 0.062 \\
 & MIDOG 2021 & 3 & 16/19 & -0.94 & 6 & 0.250 \\
 & MIDOG++ & 7 & 10/19 & -0.03 & 0 & 0.281 \\
\bottomrule
\end{tabular*}

\end{table}
